\sscp{Alienable-inalienable distinction}{13-03-05}
\citet[258f]{Blust1993} initially wanted to ascribe the development
of a distinction between \it{alienable
and inalienable} possession in Austronesian languages in our
target region to a morphosyntactic innovation in support of PCEMP.
In \srf{02-03-01-01}, we point out that the alienable-inalienable
distinction is not found in many languages in the region,
and the forms, functions,
and areas of cross-over between the two vary significantly.

We also note in \srf{02-03-01-01},
a number of scholars have shown that the alienable-inalienable
distinction is another areal development due to long-term contact
with Papuan languages, and consequently cannot
be used as a basis for linguistic subgrouping.

